% ===== PRÉAMBULE CANONIQUE — à coller VERBATIM en tête de CHAQUE .tex =====
\documentclass[11pt,a4paper]{article}
\usepackage[utf8]{inputenc}\usepackage[T1]{fontenc}\usepackage{lmodern}
\usepackage[french]{babel}
\usepackage{amsmath,amssymb,mathtools}\usepackage{icomma}
\usepackage[version=4]{mhchem}          % \ce{...} : équations & formules
\usepackage{siunitx}\sisetup{locale=FR} % \qty{}{}, \si{}, \ang{}
\usepackage{chemfig}                    % \chemfig{...} : structures développées
\usepackage{xcolor}\usepackage{colortbl}
\usepackage{tikz}\usepackage{pgfplots}\pgfplotsset{compat=1.18}
\usepackage[most]{tcolorbox}\usepackage{enumitem}\usepackage{array,booktabs}
\usepackage[a4paper,margin=2.2cm]{geometry}
\usetikzlibrary{arrows.meta,calc,angles,quotes,positioning,patterns,backgrounds,shapes.geometric}
\definecolor{primary}{RGB}{222,49,99}\definecolor{primarydark}{RGB}{150,22,55}
\definecolor{defcol}{RGB}{0,114,178}\definecolor{methcol}{RGB}{52,92,148}
\definecolor{excol}{RGB}{0,135,90}\definecolor{attncol}{RGB}{200,40,55}
\tcbset{boxbase/.style={enhanced,breakable,boxrule=0.9pt,arc=2.5pt,
  left=3mm,right=3mm,top=2.3mm,bottom=2.3mm,fonttitle=\bfseries,coltitle=white}}
% --- boîtes TUTORIEL (noms EXACTS, ne pas renommer) ---
\newtcolorbox{cle}[1][]{boxbase,colback=defcol!6,colframe=defcol,
  title={\textsc{À retenir}\ifx\relax#1\relax\relax\else\textmd{ : \itshape #1}\fi}}
\newtcolorbox{methode}[1][]{boxbase,colback=methcol!7,colframe=methcol,sharp corners,
  title={\textsc{Méthode}\ifx\relax#1\relax\relax\else\textmd{ --- \itshape #1}\fi}}
\newtcolorbox{exemplebox}[1][]{boxbase,colback=excol!5,colframe=excol,
  title={\textsc{Exemple}\ifx\relax#1\relax\relax\else\textmd{ : \itshape #1}\fi}}
\newtcolorbox{piege}[1][]{boxbase,colback=attncol!6,colframe=attncol,
  title={\textsc{Piège}\ifx\relax#1\relax\relax\else\textmd{ : \itshape #1}\fi}}
\newtcolorbox{remarque}[1][]{boxbase,colback=black!4,colframe=black!45,
  title={\textsc{Remarque}\ifx\relax#1\relax\relax\else\textmd{ : \itshape #1}\fi}}
% --- boîtes EXERCICE / CORRIGÉ (noms EXACTS, auto-comptées) ---
\newtcolorbox[auto counter]{exo}[1][]{boxbase,colback=primary!4,colframe=primary,
  title={\textsc{Exercice}~\thetcbcounter\ifx\relax#1\relax\relax\else\textmd{ --- \itshape #1}\fi}}
\newtcolorbox[auto counter]{corrige}[1][]{boxbase,colback=excol!4,colframe=excol,
  title={\textsc{Corrigé}~\thetcbcounter\ifx\relax#1\relax\relax\else\textmd{ --- \itshape #1}\fi}}
\newcommand{\R}{\mathbb{R}}\newcommand{\N}{\mathbb{N}}\newcommand{\Z}{\mathbb{Z}}\newcommand{\ds}{\displaystyle}
% (un \usepackage{fancyhdr}+en-tête est possible ; il est ignoré côté web)
% =========================================================================
\begin{document}

\begin{exo}[lire $Z$ et $A$ dans la notation]
Pour chacun de ces atomes neutres, lis la notation $^{A}_{Z}\mathrm{X}$ et donne le numéro atomique
$Z$, le nombre de masse $A$, le nombre de protons $n(p^+)$ et le nombre de nucléons.
\begin{enumerate}[label=\alph*)]
  \item $\ce{^{19}_{9}F}$
  \item $\ce{^{31}_{15}P}$
  \item $\ce{^{7}_{3}Li}$
\end{enumerate}
\end{exo}

\begin{exo}[compter les neutrons]
Le nombre de neutrons s'obtient par $n(n^0) = A - Z$. Détermine-le pour chaque noyau.
\begin{enumerate}[label=\alph*)]
  \item $\ce{^{14}_{7}N}$
  \item $\ce{^{23}_{11}Na}$
  \item $\ce{^{40}_{18}Ar}$
\end{enumerate}
\end{exo}

\begin{exo}[électrons d'un atome neutre]
Un atome neutre possède autant d'électrons que de protons. Donne le nombre d'électrons $n(e^-)$ de
chacun de ces atomes neutres.
\begin{enumerate}[label=\alph*)]
  \item $\ce{^{12}_{6}C}$
  \item $\ce{^{27}_{13}Al}$
  \item $\ce{^{32}_{16}S}$
\end{enumerate}
\end{exo}

\begin{exo}[écrire la notation d'un noyau]
On te donne la composition du noyau d'un atome neutre. Écris sa notation $^{A}_{Z}\mathrm{X}$ (avec le
symbole fourni) et donne son nombre de nucléons.
\begin{enumerate}[label=\alph*)]
  \item chlore ($\mathrm{Cl}$) : $17$ protons et $18$ neutrons ;
  \item sodium ($\mathrm{Na}$) : $11$ protons et $12$ neutrons ;
  \item oxygène ($\mathrm{O}$) : $8$ protons et $8$ neutrons.
\end{enumerate}
\end{exo}

\begin{exo}[charge d'un ensemble de particules]
La charge élémentaire vaut $e = \SI{1,6e-19}{\coulomb}$ : un proton porte $+e$, un électron porte
$-e$. Calcule la charge électrique totale (en coulombs, signe compris) de :
\begin{enumerate}[label=\alph*)]
  \item un seul proton ;
  \item un ensemble de $4$ protons ;
  \item un ensemble de $5$ électrons.
\end{enumerate}
\end{exo}

\end{document}
